%% kapa: exam=2 round=29 year=2018 period=1 subject=appraisal_practice subject_ko=감정평가실무 seq=1 number=1 points=40
%% source: Q-Net 감정평가사 기출문제 (cst00309) artlSeq=5203398 file=raw/1008/5203398/1교시 감정평가실무.pdf
%% conversion: kordoc 4.24.51 -> tools/gen2.py (auto); review=proofread-sonnet; flags=tikz
\begin{longquestion}{【문제 1】}
감정평가사 甲은 철도건설사업과 관련하여 지하공간 사용에 따른 보상목적의 감정평가를 의뢰받았다. 관련 법규 및 감정평가이론을 참작하고 제시된 자료를 활용하여 다음의 물음에 답하시오. (40점) \par
\medskip\noindent 물음 1) 감정평가사 甲은 대상토지의 지역요인을 분석하여 인근지역, 동일수급권, 유사지역의 범위를 판정하려고 한다. 인근지역의 개념과 판정기준에 대해 설명하고, 제시된 자료를 활용하여 표준지 기호 1과 기호 2, 보상선례 토지에 대해 각각 대상토지와 인근지역의 여부를 판정하시오. (10점) \par
\medskip\noindent 물음 2) 지하공간 사용에 대한 보상금을 산정하기 위한 대상토지의 적정가격을 감정평가 하시오. (15점) \par
\medskip\noindent 물음 3) 대상토지의 지하공간 사용에 대한 보상금을 산정하시오. (10점) \par
\medskip\noindent 물음 4) 관련 법규상 지하공간 사용에 대한 보상금을 감정평가하는 기준의 문제점에 대해 설명하시오. (5점) \par
< 자료 1 > 공익사업에 관한 사항 \par
1. 사업명: ○○ ～ ○○간 철도건설사업 \par
2. 사업시행자: ○○공단 \par
3. 사업추진일정 \par
1) 기본계획의 수립ㆍ고시일: 2017. 02. 02. \par
2) 보상계획 공고일: 2017. 08. 08. \par
3) 실시계획승인ㆍ고시일: 2018. 06. 06. \par
4. 권원확보방법: 구분지상권 설정 \par
< 자료 2 > 감정평가의 기본적 사항 \par
1. 대상물건: 경기도 B시 C동 산1번지의 지하터널 사용부분 \par
2. 구분지상권 설정(예정)면적: 1,200 m\textsuperscript{2} \par
3. 감정평가목적: 협의보상 \par
4. 가격시점(기준시점): 2018. 06. 01. \par
< 자료 3 > 대상토지에 관한 사항 \par
1. 소재지: 경기도 B시 C동 산1번지 \par
2. 면적: 12,000 m\textsuperscript{2}, 지목: 임야, 실제 이용상황: 자연림 \par
3. 토지이용계획: 자연녹지지역, 도시․군계획시설 공원 저촉(100 \%) \par
4. 등기사항증명서의 확인사항: 구분지상권이 설정됨 \par
1) 구분지상권자: ○○전력공사 \par
2) 목적: 154 kV 가공 송전선로 건설 \par
3) 범위: 동측 토지 상공 30 m에서 60 m 까지의 공중공간 \par
4)구분지상권 설정면적: 1,800 m\textsuperscript{2}, 존속기간: 해당 송전선로 존속시까지 \par
5) 구분지상권 설정일: 2010. 09. 09. \par
< 자료 4 > 지하공간 사용에 관한 사항 \par
1. 지하시설물의 유형: 지하터널 \par
2. 지하시설물의 크기: 높이 3 m, 너비 8 m \par
3. 토피: 대상토지가 완경사로서 위치마다 토피가 다르며, 최소 15 m ～ 최대 22 m 임(사업시행자에게 질의한 결과 평균 토피는 18 m 임) \par
4. 지하시설물 사용기간: 지하터널 존속시까지 \par
< 자료 5 > 표준지공시지가 자료 \par
{\small\centering
\begin{tabular}{|c|c|c|c|c|c|c|c|c|c|}
\hline
\multirow{2}{*}{\makecell{기호}} & \multirow{2}{*}{\makecell{소재지}} & \multirow{2}{*}{\makecell{면적\newline{}(m\textsuperscript{2})}} & \multirow{2}{*}{\makecell{지목}} & \multirow{2}{*}{\makecell{이용상황}} & \multirow{2}{*}{\makecell{용도\newline{}지역}} & \multirow{2}{*}{\makecell{도로\newline{}접면}} & \multirow{2}{*}{\makecell{형상\newline{}지세}} & \multicolumn{2}{|c|}{\makecell{공시지가(원/m\textsuperscript{2})}} \\
\cline{9-10}
 &  &  &  &  &  &  &  & \makecell{2017년} & \makecell{2018년} \\
\hline
\makecell{1} & \makecell{경기도 B시\newline{}C동 산11} & \makecell{10,000} & \makecell{임야} & \makecell{자연림} & \makecell{자연녹지} & \makecell{맹지} & \makecell{부정형\newline{}완경사} & \makecell{62,000} & \makecell{66,000} \\
\hline
\makecell{2} & \makecell{경기도 E시\newline{}F동 산20} & \makecell{12,500} & \makecell{임야} & \makecell{자연림} & \makecell{자연녹지} & \makecell{세로\newline{}(불)} & \makecell{세장형\newline{}완경사} & \makecell{58,000} & \makecell{60,000} \\
\hline
\end{tabular}\par}

※ 표준지 기호 1은 도시․군계획시설 공원에 100 \% 저촉함 \par
※ 표준지 기호 2는 도시자연공원구역에 100 \% 저촉함 \par
※ 도시․군계획시설 공원 또는 도시자연공원구역에 저촉하는 표준지의 경우 해당부분에 대해 공시지가의 감정평가시 40 \%의 감가율을 적용함 \par
※ 표준지 기호 1과 표준지 기호 2에는 154 kV 가공 송전선로 건설로 인한 구분지상권이 설정되어 있음 \par
< 자료 6 > 지가변동률 자료 \par
1. 경기도 B시 \par
{\small\centering
\begin{tabular}{|c|c|c|c|c|}
\hline
\multicolumn{2}{|c|}{\multirow{2}{*}{\makecell{구분}}} & \multicolumn{3}{|c|}{\makecell{지가변동률(단위: \%)}} \\
\cline{3-5}
\multicolumn{2}{|c|}{} & \makecell{2017년\newline{}(누계)} & \makecell{2018년\newline{}(1월 ∼ 4월 누계)} & \makecell{2018년 4월\newline{}(당월)} \\
\hline
\makecell{용도지역별} & \makecell{자연녹지} & \makecell{2.010} & \makecell{0.890} & \makecell{-0.005} \\
\hline
\makecell{이용상황별} & \makecell{임야} & \makecell{2.110} & \makecell{0.990} & \makecell{-0.003} \\
\hline
\end{tabular}\par}

※ 2018년 5월 이후 지가변동률은 미고시상태임 \par
2. 경기도 E시 \par
{\small\centering
\begin{tabular}{|c|c|c|c|c|}
\hline
\multicolumn{2}{|c|}{\multirow{2}{*}{\makecell{구분}}} & \multicolumn{3}{|c|}{\makecell{지가변동률(단위: \%)}} \\
\cline{3-5}
\multicolumn{2}{|c|}{} & \makecell{2017년\newline{}(누계)} & \makecell{2018년\newline{}(1월 ∼ 4월 누계)} & \makecell{2018년 4월\newline{}(당월)} \\
\hline
\makecell{용도지역별} & \makecell{자연녹지} & \makecell{2.120} & \makecell{1.008} & \makecell{0.002} \\
\hline
\makecell{이용상황별} & \makecell{임야} & \makecell{2.450} & \makecell{0.990} & \makecell{-0.002} \\
\hline
\end{tabular}\par}

※ 2018년 5월 이후 지가변동률은 미고시상태임 \par
< 자료 7 > 보상선례 자료 \par
1. 소재지: 경기도 E시 F동 산50번지 \par
2. 공익사업의 종류: 송전선로 건설사업(철탑부지) \par
3. 권원확보방법: 소유권 취득 \par
4. 보상액: 80,000원/m\textsuperscript{2} \par
5. 가격시점(기준시점): 2018. 05. 01. \par
6. 면적: 15,000 m\textsuperscript{2}, 지목: 임야, 실제 이용상황: 자연림 \par
7. 토지이용계획: 자연녹지지역, 도시자연공원구역(100 \%) \par
8. 도시자연공원구역에 저촉하는 토지는 보상목적의 감정평가시 40 \%의 감가율을 적 \par
용함 \par
< 자료 8 > 토지의 지역요인에 관한 자료 \par
1. 경기도 B시 C동과 E시 F동은 서로 지리적으로 접하고 있음 \par
2. 대상토지, 표준지 기호 1과 기호 2, 보상선례 토지는 서로 대체․경쟁관계가 성립하 \par
고 가격(가치)형성에 서로 영향을 미치고 있음 \par
3. 대상토지, 표준지 기호 1과 기호 2, 보상선례 토지는 모두 완경사의 국도주변 야산 지대에 속하고, 소나무와 활잡목이 혼재한 자연림지대로서 가격(가치)형성요인 중 \par
지역요인이 같거나 유사하며 가격(가치)수준이 동일함 \par
4. 대상토지와 표준지 기호 1이 속한 B시 C동과 표준지 기호 2와 보상선례 토지가 속한 E시 F동 사이에는 중앙분리대가 있는 왕복 4차선의 국도가 개설되어 있음 \par
< 자료 9 > 토지의 위치도 \par
\begin{center}
\begin{tikzpicture}[x=0.0125cm,y=-0.0125cm,font=\small,>=Stealth]
  \node[anchor=west] at (15,22) {축척없음};
  \node[font=\bfseries] at (210,72) {B시 C동};
  \node[font=\bfseries] at (470,72) {E시 F동};
  \draw[|->] (335,72) -- (305,72);
  \draw[|->] (358,72) -- (390,72);
  \draw[dotted] (342,70) -- (342,560);
  \draw (367,95) -- (278,560);
  \draw (383,95) -- (296,560);
  \node[anchor=east] at (240,218) {대상토지};
  \node at (255,218) {$\odot$};
  \node at (420,215) {$\Diamond$};
  \node[anchor=west] at (452,215) {표준지 기호 2};
  \node at (420,357) {$\square$};
  \node[anchor=west] at (452,357) {보상선례 토지};
  \node[anchor=east] at (262,424) {국도};
  \draw[->] (298,424) -- (268,424);
  \node[anchor=east] at (190,466) {표준지 기호 1};
  \node at (222,466) {$\Diamond$};
\end{tikzpicture}
\end{center}
< 자료 10 > 토지의 개별요인에 관한 자료 \par
각 토지의 개별요인에 관한 자료는 다음의 표와 같으며, 세항목별 격차율은 감정평가사가 판단할 사항임 \par
{\small\centering
\begin{tabular}{|c|c|c|c|c|}
\hline
\makecell{구분} & \makecell{대상토지} & \makecell{표준지 기호 1} & \makecell{표준지 기호 2} & \makecell{보상선례\newline{}토지} \\
\hline
\makecell{면적} & \makecell{12,000 m\textsuperscript{2}} & \makecell{10,000 m\textsuperscript{2}} & \makecell{12,500 m\textsuperscript{2}} & \makecell{15,000 m\textsuperscript{2}} \\
\hline
\makecell{취락과의 거리} & \makecell{500 m} & \makecell{750 m} & \makecell{750 m} & \makecell{750 m} \\
\hline
\makecell{도로의 상태} & \makecell{폭3m} & \makecell{맹지} & \makecell{폭3m} & \makecell{맹지} \\
\hline
\makecell{방위} & \makecell{동향} & \makecell{남향} & \makecell{동향} & \makecell{동향} \\
\hline
\makecell{경사도} & \makecell{14°} & \makecell{10°} & \makecell{10°} & \makecell{14°} \\
\hline
\makecell{형상} & \makecell{세장형} & \makecell{부정형} & \makecell{세장형} & \makecell{부정형} \\
\hline
\makecell{용도지역} & \makecell{자연녹지지역} & \makecell{자연녹지지역} & \makecell{자연녹지지역} & \makecell{자연녹지지역} \\
\hline
\makecell{도시․군\newline{}계획시설} & \makecell{공원(100 \%)} & \makecell{공원(100 \%)} & \makecell{없음} & \makecell{없음} \\
\hline
\makecell{도시자연\newline{}공원구역} & \makecell{없음} & \makecell{없음} & \makecell{저촉(100 \%)} & \makecell{저촉(100 \%)} \\
\hline
\makecell{구분지상권\newline{}(설정면적)} & \makecell{설정(1,800 m\textsuperscript{2})} & \makecell{설정(1,500 m\textsuperscript{2})} & \makecell{설정(2,500 m\textsuperscript{2})} & \makecell{없음} \\
\hline
\end{tabular}\par}

< 자료 11 > 구분지상권의 감가율 \par
154 kV 가공 송전선으로 인한 구분지상권이 설정된 임야(임지)는 주변의 시장조사 결과 필지별로 선하지의 면적비율에 따라 다음과 같이 토지가 감가되는 것으로 조사됨 \par
{\small\centering
\begin{tabular}{|c|c|c|c|}
\hline
\multirow{2}{*}{\makecell{구 분}} & \multicolumn{3}{|c|}{\makecell{선하지 면적비율}} \\
\cline{2-4}
 & \makecell{10 \% 미만} & \makecell{10 \% ∼ 20 \% 미만} & \makecell{20 \% ∼ 30 \% 미만} \\
\hline
\makecell{토지 감가율} & \makecell{15 \%} & \makecell{20 \%} & \makecell{25 \%} \\
\hline
\end{tabular}\par}

< 자료 12 > 입체이용률 배분표 \par
{\small\centering
\begin{tabular}{|c|c|c|c|}
\hline
\makecell{구분} & \makecell{저층시가지} & \makecell{주택지} & \makecell{농지․임지} \\
\hline
\makecell{건물 등 이용률(α)} & \makecell{0.75} & \makecell{0.70} & \makecell{0.80} \\
\hline
\makecell{지하부분 이용률(β)} & \makecell{0.10} & \makecell{0.15} & \makecell{0.10} \\
\hline
\makecell{그 밖의 이용률(γ)} & \makecell{0.15} & \makecell{0.15} & \makecell{0.10} \\
\hline
\makecell{(γ)의 상하배분비율} & \makecell{1:1 ∼ 3:1} & \makecell{1:1 ∼ 3:1} & \makecell{1:1 ∼ 4:1} \\
\hline
\end{tabular}\par}

※ 이용저해심도가 높은 터널 토피 20 m 이하의 경우에는 (γ)의 상하배분비율을 최고 치를 적용함 \par
< 자료 13 > 심도별 지하이용저해율 표 \par
{\small\centering
\begin{tabular}{|c|c|c|c|c|}
\hline
\makecell{한계심도} & \multicolumn{2}{|c|}{\makecell{30 m}} & \multicolumn{2}{|c|}{\makecell{20 m}} \\
\hline
\multirow{2}{*}{\makecell{체감율(\%)\newline{}토피심도(m)}} & \multirow{2}{*}{\makecell{p}} & \makecell{β × p} & \multirow{2}{*}{\makecell{p}} & \makecell{β × p} \\
\cline{3-3} \cline{5-5}
 &  & \makecell{0.15 × p} &  & \makecell{0.10 × p} \\
\hline
\makecell{0 ∼ 5 미만} & \makecell{1.000} & \makecell{0.150} & \makecell{1.000} & \makecell{0.100} \\
\hline
\makecell{5 ∼ 10 미만} & \makecell{0.833} & \makecell{0.125} & \makecell{0.750} & \makecell{0.075} \\
\hline
\makecell{10 ∼ 15 미만} & \makecell{0.667} & \makecell{0.100} & \makecell{0.500} & \makecell{0.050} \\
\hline
\makecell{15 ∼ 20 미만} & \makecell{0.500} & \makecell{0.075} & \makecell{0.250} & \makecell{0.025} \\
\hline
\makecell{20 ∼ 25 미만} & \makecell{0.333} & \makecell{0.050} & \makecell{-} & \makecell{-} \\
\hline
\end{tabular}\par}

※ p는 심도별 지하이용효율 \par
< 자료 14 > 한계심도에 관한 사항 \par
1. 한계심도는 주택지는 30 m, 농지 및 임지는 20 m임 \par
2. 한계심도를 초과하는 경우 보상율은 1.0 \% 이하임 \par
\end{longquestion}
